\section{Stable diagrams and amalgamation at closed levels}
\label{lev:section-diagrams}

Throughout this section, \(\K\) is the AEC fixed after
Proposition~\ref{found:uniform-reduction}: it is categorical in unboundedly
many cardinals, and \(\ell=\LS(\K)\). We use the order presentation \(E\)
of Theorem~\ref{found:order-template}. A tuple may have any set as its
index set. Comparisons between tuples always specify the identification
of their index sets. Reindexing by an ordinal, padding, and repetitions
do not change any of the arguments below.

The goal is to turn bounded-tuple reversal into amalgamation at selected
levels.  We first stabilize tuple comparisons across categoricity sizes,
then use the order presentation to test and realize those comparisons
over an entire base.  Amalgamation will be a conclusion of this process.

\subsection{Stabilization of strong-map classes}

Fix an infinite cardinal \(\sigma\). At a categoricity cardinal
\(\kappa>\ell+\sigma\) where Theorem~\ref{found:reversal} holds for
\(\sigma\)-tuples, put
\[
 (M,a)\equiv^\sigma_\kappa(N,b)
 \quad\Longleftrightarrow\quad
 \text{some strong map }f:M\longrightarrow N\text{ satisfies }f(a)=b.
\]
Here \(\|M\|=\|N\|=\kappa\), and the lengths of the displayed
tuples are at most \(\sigma\).
This is an equivalence relation. Reflexivity is immediate. Categoricity
allows one to regard all the maps as maps between copies of one model;
reversal then gives symmetry, and composition gives transitivity.
Using structures with universe \(\kappa\), or one fixed representative
and its tuples, its classes form a set. Their number, for all the tuple
lengths under consideration, is bounded by the fixed cardinal supplied
by Lemma~\ref{found:orbit-bound}. Denote an infinite such bound by
\(B_\sigma\). It depends on \(\ell,\sigma\), not on \(\kappa\).

\begin{lemma}[Canonical classes above one categoricity size]
\label{lev:canonical-classes}
For \(\|N\|\ge\kappa\), a tuple \(a\) in \(N\) has a well-defined
\(\equiv^\sigma_\kappa\)-class: compute it in any strong
\(\kappa\)-submodel containing \(a\). Strong embeddings between models
of size at least \(\kappa\) preserve this class.
\end{lemma}

\begin{proof}
The downward LS axiom supplies the submodel, after adding
\(\kappa\) many elements to the requested set if necessary. If
\(M_0,M_1\le_\K N\) are two such submodels, apply LS to
\(|M_0|\cup|M_1|\). It gives \(M_2\le_\K N\) of size \(\kappa\)
containing both. Coherence gives \(M_i\le_\K M_2\). The two occurrences
of \(a\) therefore have the same class. For a strong embedding
\(f:N\to N'\), the image of a witnessing submodel is a strong
\(\kappa\)-submodel of \(N'\), so the same observation proves
preservation. All these comparisons take place inside a given ambient
model; no amalgamation property is used.
\end{proof}

\begin{proposition}[Stable diagrams]
\label{lev:stabilization}
There is a definable, nondecreasing cardinal function \(t\), with
\(t(\sigma)>\ell+\sigma\) a categoricity cardinal for every infinite
\(\sigma\), having the following properties.
\begin{enumerate}
\item Tuples of length at most \(\sigma\) in models of size at least
\(t(\sigma)\) have canonical diagrams. Equality of diagrams is
preserved by every strong embedding between such models, and there
are at most \(B_\sigma\) diagrams at these lengths.
\item Diagrams include atomic diagrams, respect reindexing and
restriction, and are compatible when the bound \(\sigma\) is
increased. We denote equality by \(D(a;M)=D(b;N)\), suppressing the
bound when it is understood.
\item If \(\kappa\ge t(\sigma)\) is a categoricity cardinal, equality
of diagrams between tuples in \(\kappa\)-models is implemented by
a strong map between those models.
\item A possible diagram of a pair of tuples \((a,c)\), of total
length at most \(\sigma\), can be realized over any tuple with the
diagram of \(a\) in any categoricity model of size at least
\(t(\sigma)\). ``Possible'' means that the diagram occurs in a
model in the domain of its definition.
\end{enumerate}
\end{proposition}

\begin{proof}
If \(\gamma\ge\kappa\) are categoricity cardinals above the reversal
threshold, Lemma~\ref{lev:canonical-classes} associates a
\(\kappa\)-class to every \(\gamma\)-class. This gives a projection
\(\pi_{\gamma\kappa}\), and these projections compose. They are
surjective. Indeed, a \(\kappa\)-model is isomorphic to \(E(I)\) for
\(|I|=\kappa\); extend \(I\) to an order of size \(\gamma\) and use
the strong order map. Every tuple class in the smaller model then
occurs in the larger one. This argument supplies the needed extension
without a no-maximal-model assumption.

These projections are eventually all bijective. Otherwise, construct
a sequence of categoricity cardinals with more than \(B_\sigma\)
successive noninjective projections, and choose one final
categoricity cardinal \(\rho\) above the entire sequence. For each
noninjective projection, surjectivity supplies a pair of
\(\rho\)-classes which agree when projected to the earlier stage and
differ when projected to the later stage. Pairs selected at distinct
stages are distinct: after a pair has separated it remains separated
under every later refinement. This gives more than
\(B_\sigma\) pairs from a set of size at most \(B_\sigma\), a
contradiction. The construction uses only a set-length sequence of
cardinals and choices from sets.

Choose a categoricity cardinal \(c_\sigma\) beyond this stabilization
and reversal threshold. The canonical \(c_\sigma\)-class defines
the desired diagram on models of size at least \(c_\sigma\).
The bijective projections show that its equality can equally be
computed at every later categoricity cardinal. If two tuple bounds
are used, compare them at a categoricity cardinal above both
thresholds. Restriction of strong maps then gives the same
restricted equivalence. Atomic formulas and their negations are
preserved because strong maps are structure embeddings.

For definiteness, thresholds can be chosen as least cardinals having
the stated eventual properties. Choose \(t(\sigma)\) to be the
least categoricity cardinal strictly above \(\ell+\sigma\) and
above all the thresholds \(c_\tau\) for infinite \(\tau\le\sigma\).
This is a nondecreasing definable function. The preceding arguments
prove the first three assertions.

For the fourth, first put a realization of \((a,c)\) into a strong
\(c_\sigma\)-submodel. If the originally chosen ambient size is
smaller than a desired target categoricity size, this categorical
submodel can be extended to that size by the order presentation.
If the original ambient size is larger, LS gives a submodel of the
desired size containing the realization. Thus the joint diagram
has a realization in a model of the target categoricity size.
Apply the third assertion to the first tuple and the specified tuple
in the target. The resulting strong map transports the whole
realization as required.
\end{proof}

\begin{proposition}[Categoricity at countably cofinal closure points]
\label{lev:closure-categoricity}
Suppose \(\zeta>\ell\) is an infinite cardinal of countable
cofinality and
\[
 t(\sigma)<\zeta\qquad\text{for every infinite }\sigma<\zeta.
\]
Then \(\K\) is categorical in \(\zeta\). Moreover any diagram match
between bounded tuples in two \(\zeta\)-models extends to an
isomorphism between those models. Here and throughout the paper,
bounded at a level means of cardinality strictly less than that level.
\end{proposition}

\begin{proof}
Existence follows from \(E(I)\) for \(|I|=\zeta\). Let \(M,N\)
have size \(\zeta\), and choose countable increasing exhaustions of
their underlying sets by sets of size less than \(\zeta\).
Start with the prescribed bounded diagram match, or with the empty
match. At a forward step, enumerate together the current source
tuple and the next part of the exhaustion of \(M\). Its length and
the length of the current target tuple are bounded by an infinite
\(\sigma<\zeta\). Put the data on each side inside a strong
submodel of size \(t(\sigma)<\zeta\). Their current tuples have
the same diagram. Proposition~\ref{lev:stabilization} gives a
strong map of these submodels matching those tuples. Its values
on the new entries extend the match, and the enlarged tuples
again have equal diagrams. At the next step interchange \(M,N\).

Only finitely many previously chosen bounded tuples are used at
each step, so their combined length stays below \(\zeta\).
The union of the countable sequence of partial matches is a
bijection \(|M|\to|N|\). Every finite tuple and every value of a
finitary function occur at some step, where the atomic diagram
is preserved. The bijection is therefore a structure isomorphism.
\end{proof}

The closure cardinals in Proposition~\ref{lev:closure-categoricity}
are unbounded. For example, starting above any specified cardinal,
choose a strictly increasing sequence \((\alpha_n:n<\omega)\)
such that \(\alpha_{n+1}\) exceeds all \(t(\sigma)\) for
\(\sigma\le\alpha_n\). Its supremum has the required closure and
countable cofinality. Including \(2^{\alpha_n}\) in these requirements
makes the supremum a strong limit. We will use these stronger
closure points below.

\subsection{Natural diagrams and interchange of blocks}

\begin{definition}[Natural and placed diagrams]
\label{lev:natural-diagrams}
Consider a set-sized list of terms in an order presentation. Each
term consists of a label and an increasing finite list of positions.
Extend the order to size at least the threshold for the list and
compute its diagram there. The result is its \emph{natural diagram}.
If \(A=E(I)\), a \emph{placement} of an additional tuple \(d\)
over \(A\) is a labeled representation in an order extension of
\(I\), retaining the positions and chosen representations of the
parameters in \(A\). Its \emph{placed comparison over \(A\)} is
the natural diagram of \((d,a_A)\), where \(a_A\) is a fixed
enumeration of \(A\).
\end{definition}

The natural diagram is independent of the extension used. Indeed,
linear orders amalgamate over a common suborder: for each cut of
the suborder, place the remaining portions of the two orders in
that cut while preserving each portion's order. Extending the
amalgam further if necessary gives a common sufficiently large
presentation. Strong order maps and
Proposition~\ref{lev:stabilization} identify the two diagrams.
The same argument shows that only the labels and the joint
order/equality pattern of their supports matter. All diagrams
with parameters use matching parameter enumerations. One may
instead list all labels on the base order, allowing repetitions:
restriction gives an enumeration of the elements, and reinserting
the repetitions is forced by their equalities.

A placed comparison over a whole \(\mu\)-model is evaluated in a
sufficiently large presentation. It is not, by this definition,
an assertion that the original \(\mu\)-model itself lies above the
threshold for diagrams of length \(\mu\).

\begin{lemma}[Orbits over parameters]
\label{lev:parameter-orbits}
Let \(\sigma\) be infinite and suppose
\[
 \chi^\sigma=\chi,
 \qquad \chi\ge \ell+2^{2^\sigma}.
\]
At every sufficiently large categoricity cardinal, there are at
most \(\chi\) automorphism orbits of \(\sigma\)-tuples over any
specified set of at most \(\chi\) parameters.
\end{lemma}

\begin{proof}
Present the categoricity model as \(E(H(\delta))\), choosing
\(|\delta|\) to be its cardinality and \(\cf(\delta)>\chi\).
This is possible at every sufficiently large cardinal, including
singular cardinals. Fix term representations of the parameters,
and let \(S\subseteq H(\delta)\) be their support set. Its size
is at most \(\chi\). Every member of \(S\) is coded by a finite
list of ordinals and integer and rational coefficients. The union
of its ordinal supports consequently also has size at most
\(\chi\).

For one further position, its finite ordinal support has at most
\(\chi\) possible order/equality patterns over these old ordinal
supports. The discrete coefficients contribute only countably
many choices for each finite code. Such a pattern determines
all comparisons of the new position with \(S\). For at most
\(\sigma\) new positions, include their mutual comparisons,
their identifications with old positions, and the term labels.
There are at most \(\chi^\sigma=\chi\) choices so far.

To apply Theorem~\ref{found:order-orbits}, also record which strict
slots over the old and new positions are realized in
\(H(\delta)\). The old support \(S\) is fixed. A slot over \(S\)
not containing a new position is unchanged. At most \(\sigma\)
old slots contain new positions, and inside each there are at
most \(2^\sigma\) cuts of the new support to test. The new flags
therefore have at most \(2^{2^\sigma}\le\chi\) possible values.
The full record has at most \(\chi\) possibilities.

The combined list of old and new positions has size at most
\(\chi<\cf(\delta)\). Theorem~\ref{found:order-orbits} extends
equality of full records to an order automorphism fixing \(S\)
pointwise. Its induced automorphism of the presentation fixes
the parameters and matches the tuples. Thus the number of
structure-automorphism orbits is at most the number of records.
\end{proof}

\begin{lemma}[Disjoint block transpositions]
\label{lev:block-transpositions}
Take two set-sized families of labeled terms. In disjoint convex
regions of their joint support order, suppose each region consists
of two adjacent blocks, one used only by the first family and one
only by the second. Interchange the blocks in every such region,
retaining their internal orders and all positions outside the
regions. The joint natural diagram is unchanged.
\end{lemma}

\begin{proof}
Suppose the original and modified diagrams were different. Let
\(\sigma\) bound the two lists, their supports, and the shared
positions. Choose \(\chi\ge\ell+2^{2^\sigma}\) with
\(\chi^\sigma=\chi\), for instance by taking a sufficiently
large power of two.

In each region repeat the pair of blocks along an arbitrary
linear order \(L\). If the first family's block was on the left,
use the order \(L\) for the repetitions; if it was on the right,
use the reverse of \(L\). Write \(d_i\) for the first-family
tuple using the blocks at index \(i\), and \(e_j\) for the
second-family tuple using the blocks at index \(j\); retain all
the shared positions outside the regions. For \(i<j\), the
diagram of \((d_i,e_j)\) is the original one, and for \(i>j\)
it is the modified one. Comparisons within either family and
across distinct regions do not change.

Let \(\epsilon\) be the least infinite cardinal satisfying
\(2^\epsilon>\chi\). The binary tree \(2^{<\epsilon}\) has
at most \(\chi\) nodes: for \(\alpha<\epsilon\) each level has
size at most \(\chi\), and \(\epsilon\le\chi\). Give it
lexicographically ordered divider points, one between the two
successor subtrees at each node. Explicitly, encode branch
digits by \(-1,1\), and encode the divider at a node \(s\)
by the sequence \(s\) followed by \(0\), using
\(-1<0<1\) in lexicographic comparison. Adjoin one point for every
branch in \(2^\epsilon\). Distinct branch points determine
distinct strict cuts over the divider points. Use this as the
order \(L\).

The tuples \(e_j\) at the divider indices, together with the
fixed shared data, use at most \(\chi\) parameters. If two
branch-indexed tuples \(d_i\) were in the same automorphism
orbit over these parameters, their diagrams with every divider
tuple would agree. A divider between the two branches gives
the two different diagrams, a contradiction. Hence there are
more than \(\chi\) such orbits. Pad the constructed order to a
sufficiently large categoricity size and apply
Lemma~\ref{lev:parameter-orbits}. This is impossible.
\end{proof}

\subsection{Bounded tests over a well-ordered presentation}

The next elementary set-theoretic observation will also be used for
the more elaborate layouts of later sections.

\begin{lemma}[Ordinal-vector witnesses]
\label{lev:ordinal-vectors}
Let \(C\) be a set, and let \(U\) be a definable class of functions
\(C\to\Ord\), with fixed set parameters allowed in its definition.
There is a set \(U_0\subseteq U\) such that each \(u\in U\)
dominates some \(v\in U_0\) coordinatewise.
\end{lemma}

\begin{proof}
The assertion is immediate if \(C\) is empty or \(U\) is empty.
Well-order \(C\), and suppose no such set exists. For every
ordinal \(\beta\), the set \(U\cap V_\beta\) fails to generate
\(U\) upwards. There is therefore an element of \(U\) which
dominates none of its members. Let \(h(\beta)>\beta\) be the
least ordinal such that one such witness lies in \(V_{h(\beta)}\).
This is a definable ordinal function.

Take an ordinal length \(\lambda\) satisfying the partition
relation for pairs, \(|C|\) colors, and an infinite homogeneous
set, supplied by Lemma~\ref{found:partition}. Define increasing
ordinals \(\beta_i\), for \(i\le\lambda\), recursively, with
\(\beta_{i+1}>h(\beta_i)\) and with limit values above all earlier
values. Choice for a set-indexed collection of nonempty subsets
of \(V_{\beta_\lambda}\) now selects witnesses \(u_i\) which
dominate no member of \(U\cap V_{\beta_i}\), for \(i<\lambda\).
For \(i<j\), we have \(u_i\in V_{\beta_j}\), so some coordinate
satisfies \(u_i(c)>u_j(c)\). Color the pair \(i<j\) by the first
such coordinate in the fixed well-order of \(C\). An infinite
homogeneous set yields an infinite strictly decreasing sequence
of ordinals in that coordinate. This contradiction proves the
claim. The rank bounding above ensures that only ordinary
set-sized Choice, not a choice principle for proper classes,
was used.
\end{proof}

\begin{proposition}[Uniform testing of placed comparisons]
\label{lev:testing}
There is a definable nondecreasing infinite cardinal function
\(g\), depending on \(\K\) and its fixed presentation, such that
the following holds. If \(I\) is an infinite well-order, \(A=E(I)\), and
two placements of a tuple of length at most \(\sigma\) have
different placed comparisons over \(A\), then their natural
diagrams already differ over a subset of \(A\) of size at most
\(g(\sigma)\). The bound is independent of \(|I|\).
\end{proposition}

\begin{proof}
Amalgamate the two support orders over \(I\), and let \(T\)
be the union of the supports of the two added tuples. Its size
is at most \(\sigma\). Freeze its order, the labels and supports
of the two tuples, and the subset of its positions which belong
to \(I\). Up to a relabeling on a fixed set of size at most
\(\sigma\), these frozen data range over a set.

For fixed frozen data, each component of \(I\setminus T\) in
a strict slot over \(T\) is well-ordered. Record its ordinal
length. The set of slots is fixed, so these lengths form an
ordinal vector. Restrict attention to the vectors that give
valid configurations and for which the two placed comparisons
are unequal. This is a definable class: order presentations
and diagrams are defined from set codes, and diagrams of any
specified set length can be evaluated above their thresholds.

Coordinatewise domination gives an order embedding of the
smaller configuration into the larger one, fixing \(T\) and
preserving the base positions. Natural diagrams respect these
embeddings. Thus a valid smaller unequal configuration remains
a witness in any dominating valid configuration. Apply
Lemma~\ref{lev:ordinal-vectors} to this class, obtaining a set
of valid witnesses. Equivalently one could apply it to the
upward closure and select its generators from valid
configurations; no assertion that arbitrary vectors themselves
describe well-orders is required.

Bound the sizes of the base orders in this set of witnesses.
An embedded witness uses only that bounded number of old
positions. All labels on these positions give at most their
cardinality plus \(\ell\) elements of \(A\), and the diagrams
of the two added tuples with these elements are unequal.
Taking a cardinal bound over the set of frozen data gives
\(g(\sigma)\). Increasing it to an infinite nondecreasing
function causes no change in the assertion. When a witness
uses only finitely many positions, ``all labels'' simply means
their evaluated terms; no assertion about membership of a
finite-order presentation in \(\K\) is needed.
\end{proof}

\begin{proposition}[Compact placement]
\label{lev:compact-placement}
Let \(A=E(I)\), where \(I\) is an infinite well-order, and let
\(\chi>\ell+\sigma\) satisfy \(\chi^\sigma=\chi\). Suppose that
for every \(X\subseteq A\) of size at most \(\chi\) a possible
joint diagram \(q_X(d,X)\) is prescribed, where \(d\) has length
at most \(\sigma\). Assume these prescriptions restrict
coherently and that their parameter diagrams are the natural
ones of \(X\). Then one placement of \(d\) over \(I\) realizes
all the prescriptions.
\end{proposition}

\begin{proof}
Fix one term representation per parameter in \(A\). Write
\(S_X\subseteq I\) for the supports of \(X\). For one set
\(X\), its prescription has a placement. To see this, extend
\(I\) to a sufficiently large categoricity size and use
Proposition~\ref{lev:stabilization} to realize the possible
joint diagram over the actual parameter tuple. The realization
has term representations in that presentation.

A support profile over \(X\) records the labels for \(d\),
its formal support order \(T_0\) of size at most \(\sigma\),
and comparisons and coincidences with \(S_X\). A well-order
of size at most \(\chi\) has at most \(\chi\) cuts. Hence
there are at most \(\chi^\sigma=\chi\) profiles. The positions
of \(T_0\) in one strict slot over \(S_X\) form a cluster.
Each profile has a continuation over \(I\) which leaves
every such cluster intact: insert the cluster at a single
cut of \(I\) extending its specified cut over \(S_X\).

Call a hit of a new parameter support on \(T_0\), or the
separation of two positions in one cluster by a new parameter
support, a new event. For \(X'\supseteq X\), all continuations
of a profile with no new event over \(S_{X'}\) have the same
natural diagram over \(X'\). Within each old slot, the whole
cluster can only change its position across a block of new
parameter positions. The slots are disjoint, so
Lemma~\ref{lev:block-transpositions} applies simultaneously.

Suppose no placement solves all the prescriptions. Start with
any \(X_0\) of size at most \(\chi\). Given \(X_i\), choose
for every profile over \(X_i\) a continuation with no new
event over all of \(I\). It fails a prescription on some set
of at most \(\chi\) parameters. Adjoin one such failing set
for each profile to obtain \(X_{i+1}\). At limits take unions,
and carry out \(\sigma^+\) stages. Every set constructed,
including \(X_* = \bigcup_{i<\sigma^+}X_i\), has size at most
\(\chi\), because \(\sigma^+\le\chi\).

Take a placement realizing \(q_{X_*}\). At each stage its
profile over \(X_i\) was considered. If it had no new event
between \(S_{X_i}\) and \(S_{X_{i+1}}\), the preceding
no-event comparison would identify its diagram on the
selected failing set with that of the chosen continuation.
This contradicts its realization of \(q_{X_*}\). Thus it
has a new event at every stage. A support position can be
newly hit only once, and a pair can be newly separated only
once. There are at most \(\sigma+\sigma^2=\sigma\) such
events. They cannot occur at \(\sigma^+\) different stages.
\end{proof}

\begin{lemma}[Unique lifts of bounded diagrams]
\label{lev:placed-lifts}
Let \(\mu>\ell\) be a strong-limit cardinal of countable
cofinality strictly closed below itself under \(t\) and \(g\).
Let \(A\le_\K B\) have size \(\mu\), and identify
\(A=E(I)\) for a well-order \(I\) of size \(\mu\). For a
tuple \(d\) in \(B\) of length at most \(\sigma<\mu\), the
diagrams
\[
 D(d,X;B),\qquad X\subseteq A,\quad |X|<\mu,
\]
lift to one placed comparison over all of \(A\), uniquely
as a comparison diagram. Restricting a lifted tuple gives
the lift of its restricted diagram.
\end{lemma}

\begin{proof}
The displayed bounded diagrams are defined, since their
thresholds lie below \(\mu\), and they agree on the parameters
with the natural diagrams. For every sufficiently large
\(\chi<\mu\) with \(\chi^\sigma=\chi\),
Proposition~\ref{lev:compact-placement} gives a placement
realizing the tests of size at most \(\chi\). Such \(\chi\)
are cofinal below \(\mu\): use \(\chi=2^\alpha\) for
\(\sigma\le\alpha<\mu\), enlarging \(\alpha\) as necessary.

Once \(\chi\ge g(\sigma)\), any two resulting placements
have the same full comparison by Proposition~\ref{lev:testing}.
Fix one of them. Every bounded test is covered by a larger
choice of \(\chi\), whose placement has that same full
comparison. The fixed placement consequently realizes every
bounded test. The same argument gives uniqueness. A
restriction realizes the restricted tests and hence, by
uniqueness, is their lift.
\end{proof}

\subsection{Amalgamation and type counts}

\begin{theorem}[Amalgamation at closed levels]
\label{lev:amalgamation}
There is an unbounded class of strong-limit cardinals \(\mu\)
of countable cofinality with the following properties.
\begin{enumerate}
\item \(\K\) is categorical in \(\mu\), and any two
\(\mu\)-sized strong extensions of a \(\mu\)-model have a
strong amalgam of size \(\mu\).
\item If \(A\le_\K B_0,B_1\) all have size \(\mu\), and
\(d_i\in B_i\) are tuples of the same length at most
\(\sigma<\mu\), there is such an amalgam identifying
\(d_0\) with \(d_1\) if and only if
\[
 D(d_0,X;B_0)=D(d_1,X;B_1)
 \quad\text{for every }X\subseteq A,
 \quad |X|\le g(\sigma).
\]
\item There are at most \(\mu\) singleton types over each
\(\mu\)-model, where type means identification in an amalgam
over that model.
\end{enumerate}
One may take every sufficiently large strong-limit cardinal
of countable cofinality strictly closed under \(t\) and \(g\).
\end{theorem}

\begin{proof}
Such levels are unbounded by the countable closure construction
following Proposition~\ref{lev:closure-categoricity}.
Categoricity at them has already been proved. Identify a
given base with \(A=E(I)\), where \(I\) is a well-order of
size \(\mu\). Extend \(I\) to an order \(J\) whose cardinality
\(\gamma\) is a categoricity cardinal at least \(t(\mu)\),
and let \(G=E(J)\). Thus \(A\) is strongly included in its
specified position in \(G\).

We show that every \(\mu\)-sized strong extension \(B\) of
\(A\) embeds strongly into \(G\) over \(A\). Choose
\(M_n\le_\K B\), increasing with union \(B\), with
\(\|M_n\|<\mu\). This follows by LS from a countable
exhaustion of \(B\), containing the preceding submodel at
each step; coherence makes the chain strong. Choose compatible
enumerations \(d_n\) of these models, so that each earlier
list is a sublist of the next.

Lemma~\ref{lev:placed-lifts} assigns coherent placed comparisons
to \((d_n,A)\). Each is realized over \(A\) in \(G\), and
the realizations can retain the previous assignments. More
explicitly, realize the corresponding term pattern in an
order extension of size \(\gamma\). Its parameter tuple
\(A\), and then its tuple \((A,d_{n-1})\), have the same
diagram as the already assigned ones. Their length is at
most \(\mu\), so Proposition~\ref{lev:stabilization} at
\(\gamma\ge t(\mu)\) supplies the strong map realizing the
new list while fixing these previous data.

It remains to verify strongness of the copied small models;
atomic agreement by itself would not suffice. Fix \(n\),
and bound the length of \(d_n\) by an infinite
\(\sigma_n<\mu\). Choose the categoricity cardinal
\(\rho=t(\sigma_n)<\mu\). LS puts the original enumeration
inside \(B_n'\le_\K B\) of size \(\rho\), and the copied
enumeration inside \(G_n'\le_\K G\) of that size. Coherence
gives \(M_n\le_\K B_n'\). The copied and original lists
have the same bounded diagram; Proposition~\ref{lev:stabilization}
gives a strong map \(B_n'\to G_n'\) with exactly the prescribed
values on \(d_n\). Its restriction shows that the copied
\(M_n\) is a strong submodel of \(G\).

The copied models are nested, and coherence makes their
inclusions strong. Their union is a strong submodel of \(G\)
by the chain and common-upper-bound axioms. It is the image
of an embedding of \(B\). Every element of \(A\) eventually
occurs in some \(d_n\), and its equality with that parameter
is part of the lifted comparison. The embedding therefore
fixes \(A\). Performing this construction for \(B_0,B_1\)
in the same \(G\) gives an amalgam. LS supplies a strong
\(\mu\)-submodel of \(G\) containing both images; coherence
makes it a strong amalgam of the required size.

For the second assertion, any identifying amalgam preserves
the displayed bounded diagrams, giving necessity. Conversely,
Proposition~\ref{lev:testing} identifies the placed lifts
of \(d_0,d_1\) over \(A\). Realize this common lift once
in \(G\). On each side start the preceding countable
construction with the specified tuple already assigned,
and choose every small-model enumeration to contain that
tuple. Homogeneity over \(A\) and the old assignments works
exactly as above. The two resulting embeddings identify the
tuples, and the same LS shrinking applies.

For a singleton, a placement uses one label on a finite
support. Relative to the well-order \(I\), each support
position is either an old position or lies in one of its
\(|I|+1\) cuts. Together with the finitely many mutual
comparisons and the label, these data have at most
\(\mu\) possibilities. They determine the full natural
comparison. The second assertion identifies exactly the
types with equal comparison, proving the bound.
\end{proof}

Amalgamation has thus been obtained at the stated levels from the
AEC axioms and unbounded categoricity. We have not assumed that
all models of intermediate cardinalities amalgamate. The next
section records precisely the level-by-level properties that
will be required of later categories of finite systems.
